\textbf{Phần 1. Trắc nghiệm lựa chọn}
\setcounterpageref{ex}{0}

\begin{ex} %[5H12V71]
Một hình tam giác có cạnh đáy dài $24\text{ cm}$. Nếu kéo dài cạnh đáy thêm $6\text{ cm}$ về cùng một phía thì diện tích tam giác tăng thêm $45\text{ cm}^2$. Diện tích tam giác ban đầu là
\choice
{$150\text{ cm}^2$}
{\True $180\text{ cm}^2$}
{$210\text{ cm}^2$}
{$360\text{ cm}^2$}
\loigiai{
Chiều cao của tam giác là:
$$45 \times 2 : 6 = 15\text{ (cm)}$$
Diện tích tam giác ban đầu là:
$$\dfrac{24 \times 15}{2} = 180\text{ (cm}^2\text{)}$$
\textbf{Đáp án của câu hỏi là B}
}
\end{ex}

\begin{ex} %[5H12V71]
Một mảnh đất hình thang có đáy lớn $32\text{ m}$, đáy bé bằng $\dfrac{3}{4}$ đáy lớn, chiều cao kém đáy bé $6\text{ m}$. Diện tích của mảnh đất đó là
\choice
{$480\text{ m}^2$}
{\True $504\text{ m}^2$}
{$1\,008\text{ m}^2$}
{$576\text{ m}^2$}
\loigiai{
Độ dài đáy bé là:
$$32 \times \dfrac{3}{4} = 24\text{ (m)}$$
Chiều cao của mảnh đất hình thang là:
$$24 - 6 = 18\text{ (m)}$$
Diện tích mảnh đất là:
$$\dfrac{(32 + 24) \times 18}{2} = 504\text{ (m}^2\text{)}$$
\textbf{Đáp án của câu hỏi là B}
}
\end{ex}

\begin{ex} %[5H12V71]
Một biển báo giao thông hình tròn có đường kính $50\text{ cm}$. Người ta sơn một vành viền màu đỏ rộng $5\text{ cm}$ chạy dọc theo mép ngoài của biển báo. Diện tích phần vành viền màu đỏ là (lấy $\pi = 3{,}14$)
\choice
{$628\text{ cm}^2$}
{\True $706{,}5\text{ cm}^2$}
{$785\text{ cm}^2$}
{$1\,413\text{ cm}^2$}
\loigiai{
Bán kính hình tròn ngoài là: $50 : 2 = 25\text{ (cm)}$.\\
Bán kính hình tròn trong là: $25 - 5 = 20\text{ (cm)}$.\\
Diện tích hình tròn ngoài:
$$3{,}14 \times 25 \times 25 = 1\,962{,}5\text{ (cm}^2\text{)}$$
Diện tích hình tròn trong:
$$3{,}14 \times 20 \times 20 = 1\,256\text{ (cm}^2\text{)}$$
Diện tích phần vành viền đỏ là:
$$1\,962{,}5 - 1\,256 = 706{,}5\text{ (cm}^2\text{)}$$
\textbf{Đáp án của câu hỏi là B}
}
\end{ex}

\begin{ex} %[5H12V71]
Bánh xe của một chiếc xe đạp có bán kính $0{,}35\text{ m}$. Khi bánh xe lăn trên đường được $500$ vòng thì chiếc xe đạp đi được quãng đường dài bao nhiêu mét?
\choice
{$1\,050\text{ m}$}
{\True $1\,099\text{ m}$}
{$2\,198\text{ m}$}
{$549{,}5\text{ m}$}
\loigiai{
Chu vi bánh xe đạp là:
$$0{,}35 \times 2 \times 3{,}14 = 2{,}198\text{ (m)}$$
Quãng đường xe đạp đi được sau 500 vòng là:
$$2{,}198 \times 500 = 1\,099\text{ (m)}$$
\textbf{Đáp án của câu hỏi là B}
}
\end{ex}

\begin{ex} %[5H12V71]
Một phòng học dạng hình hộp chữ nhật có chiều dài $8\text{ m}$, chiều rộng $6\text{ m}$ và chiều cao $3{,}5\text{ m}$. Người ta sơn bốn bức tường bên trong và trần nhà. Biết tổng diện tích các cửa là $12\text{ m}^2$. Diện tích cần quét sơn là
\choice
{$146\text{ m}^2$}
{\True $134\text{ m}^2$}
{$122\text{ m}^2$}
{$98\text{ m}^2$}
\loigiai{
Diện tích xung quanh phòng học là:
$$(8 + 6) \times 2 \times 3{,}5 = 98\text{ (m}^2\text{)}$$
Diện tích trần nhà là:
$$8 \times 6 = 48\text{ (m}^2\text{)}$$
Diện tích cần quét sơn là:
$$98 + 48 - 12 = 134\text{ (m}^2\text{)}$$
\textbf{Đáp án của câu hỏi là B}
}
\end{ex}

\begin{ex} %[5H12V71]
Người ta dán giấy màu lên toàn bộ $4$ mặt xung quanh và $1$ mặt đáy của một chiếc hộp quà dạng hình lập phương không có nắp cạnh $15\text{ cm}$. Diện tích giấy màu cần dùng ít nhất là
\choice
{$900\text{ cm}^2$}
{\True $1\,125\text{ cm}^2$}
{$1\,350\text{ cm}^2$}
{$225\text{ cm}^2$}
\loigiai{
Hộp không có nắp nên diện tích cần dán gồm 5 mặt.\\
Diện tích 1 mặt của hộp là:
$$15 \times 15 = 225\text{ (cm}^2\text{)}$$
Diện tích giấy màu cần dán là:
$$225 \times 5 = 1\,125\text{ (cm}^2\text{)}$$
\textbf{Đáp án của câu hỏi là B}
}
\end{ex}

\begin{ex} %[5H12V71]
Một mảnh đất gồm một hình chữ nhật có chiều dài $16\text{ m}$, chiều rộng $10\text{ m}$ và một hình tam giác có cạnh đáy trùng với chiều rộng của hình chữ nhật, chiều cao ứng với cạnh đáy đó là $6\text{ m}$. Diện tích toàn bộ mảnh đất là
\choice
{$175\text{ m}^2$}
{\True $190\text{ m}^2$}
{$220\text{ m}^2$}
{$250\text{ m}^2$}
\loigiai{
Diện tích phần đất hình chữ nhật là:
$$16 \times 10 = 160\text{ (m}^2\text{)}$$
Diện tích phần đất hình tam giác là:
$$\dfrac{10 \times 6}{2} = 30\text{ (m}^2\text{)}$$
Diện tích toàn bộ mảnh đất là:
$$160 + 30 = 190\text{ (m}^2\text{)}$$
\textbf{Đáp án của câu hỏi là B}
}
\end{ex}

\begin{ex} %[5H12V71]
Một bể cá dạng hình hộp chữ nhật dài $60\text{ cm}$, rộng $40\text{ cm}$, cao $35\text{ cm}$. Mực nước ban đầu cao $20\text{ cm}$. Khi thả một hòn non bộ chìm hoàn toàn vào bể thì nước dâng lên thêm $5\text{ cm}$ (nước không tràn). Thể tích của hòn non bộ là
\choice
{$10\text{ dm}^3$}
{\True $12\text{ dm}^3$}
{$15\text{ dm}^3$}
{$24\text{ dm}^3$}
\loigiai{
Thể tích hòn non bộ bằng thể tích phần nước dâng lên:
$$60 \times 40 \times 5 = 12\,000\text{ (cm}^3\text{)} = 12\text{ (dm}^3\text{)}$$
\textbf{Đáp án của câu hỏi là B}
}
\end{ex}

\begin{ex} %[5H12V71]
Một thửa ruộng hình thang có chiều cao $15\text{ m}$. Nếu mở rộng đáy lớn thêm $4\text{ m}$ và đáy bé thêm $2\text{ m}$ mà vẫn giữ nguyên chiều cao thì diện tích thửa ruộng tăng thêm bao nhiêu mét vuông?
\choice
{$30\text{ m}^2$}
{\True $45\text{ m}^2$}
{$60\text{ m}^2$}
{$90\text{ m}^2$}
\loigiai{
Tổng độ dài hai đáy tăng thêm là:
$$4 + 2 = 6\text{ (m)}$$
Diện tích thửa ruộng tăng thêm là:
$$\dfrac{6 \times 15}{2} = 45\text{ (m}^2\text{)}$$
\textbf{Đáp án của câu hỏi là B}
}
\end{ex}

\begin{ex} %[5H12V71]
Một khối gỗ hình lập phương có cạnh dài $8\text{ cm}$. Người ta xẻ khối gỗ đó thành các khối gỗ nhỏ hình lập phương có cạnh dài $2\text{ cm}$. Số khối gỗ nhỏ thu được là
\choice
{$32\text{ khối}$}
{\True $64\text{ khối}$}
{$16\text{ khối}$}
{$128\text{ khối}$}
\loigiai{
Thể tích khối gỗ ban đầu là:
$$8 \times 8 \times 8 = 512\text{ (cm}^3\text{)}$$
Thể tích một khối gỗ nhỏ là:
$$2 \times 2 \times 2 = 8\text{ (cm}^3\text{)}$$
Số khối gỗ nhỏ cắt được là:
$$512 : 8 = 64\text{ (khối)}$$
\textbf{Đáp án của câu hỏi là B}
}
\end{ex}

\begin{ex} %[5H12V71]
Cho tam giác $ABC$ vuông tại $A$ có cạnh $AB = 30\text{ cm}$, cạnh $AC = 40\text{ cm}$, cạnh huyền $BC = 50\text{ cm}$. Chiều cao $AH$ hạ từ đỉnh $A$ xuống cạnh huyền $BC$ có độ dài là
\choice
{$20\text{ cm}$}
{\True $24\text{ cm}$}
{$25\text{ cm}$}
{$28\text{ cm}$}
\loigiai{
Diện tích tam giác vuông $ABC$ là:
$$\dfrac{30 \times 40}{2} = 600\text{ (cm}^2\text{)}$$
Chiều cao $AH$ có độ dài là:
$$600 \times 2 : 50 = 24\text{ (cm)}$$
\textbf{Đáp án của câu hỏi là B}
}
\end{ex}

\begin{ex} %[5H12V71]
Một bể nước hình lập phương có cạnh bên trong là $1{,}4\text{ m}$. Hiện bể đang chứa $1\,372\text{ l}$ nước. Mực nước hiện tại trong bể cao bao nhiêu xăng-ti-mét?
\choice
{$60\text{ cm}$}
{\True $70\text{ cm}$}
{$80\text{ cm}$}
{$75\text{ cm}$}
\loigiai{
Đổi: $1\,372\text{ l} = 1\,372\text{ dm}^3 = 1{,}372\text{ m}^3$.\\
Diện tích đáy bể là:
$$1{,}4 \times 1{,}4 = 1{,}96\text{ (m}^2\text{)}$$
Chiều cao của mực nước trong bể là:
$$1{,}372 : 1{,}96 = 0{,}7\text{ (m)} = 70\text{ (cm)}$$
\textbf{Đáp án của câu hỏi là B}
}
\end{ex}

\textbf{Phần 2. Trắc nghiệm đúng sai}
\setcounterpageref{ex}{0}

\begin{ex} %[5H12V71]
Một mảnh đất hình thang có đáy lớn $28\text{ m}$, đáy bé $18\text{ m}$ và chiều cao $15\text{ m}$. Người ta mở rộng mảnh đất bằng cách kéo dài đáy lớn thêm $5\text{ m}$ về một phía và giữ nguyên đáy bé cùng chiều cao. Xét tính đúng sai của các khẳng định sau:
\choiceTF
{\True Diện tích mảnh đất hình thang ban đầu là $345\text{ m}^2$.}
{Phần diện tích đất tăng thêm là một hình tam giác có diện tích $75\text{ m}^2$.}
{\True Phần diện tích đất tăng thêm là một hình tam giác có diện tích $37{,}5\text{ m}^2$.}
{\True Diện tích mảnh đất sau khi mở rộng là $382{,}5\text{ m}^2$.}
\loigiai{
\begin{itemchoice}
    \itemch Đúng. Diện tích ban đầu: $\dfrac{(28 + 18) \times 15}{2} = 345\text{ m}^2$.
    \itemch Sai. Phần tăng thêm có diện tích là $\dfrac{5 \times 15}{2} = 37{,}5\text{ m}^2 \neq 75\text{ m}^2$.
    \itemch Đúng. Diện tích tăng thêm là $37{,}5\text{ m}^2$.
    \itemch Đúng. Diện tích sau mở rộng: $345 + 37{,}5 = 382{,}5\text{ m}^2$.
\end{itemchoice}
}
\end{ex}

\begin{ex} %[5H12V71]
Một bể chứa nước dạng hình hộp chữ nhật có kích thước bên trong lòng bể: chiều dài $2{,}5\text{ m}$, chiều rộng $1{,}6\text{ m}$, chiều cao $1{,}2\text{ m}$. Xét tính đúng sai của các khẳng định sau:
\choiceTF
{\True Diện tích đáy của bể nước là $4\text{ m}^2$.}
{\True Dung tích tối đa của bể là $4{,}8\text{ m}^3$, tương đương $4\,800\text{ l}$.}
{Khi trong bể chứa $3\,200\text{ l}$ nước thì mực nước cao $0{,}75\text{ m}$.}
{\True Khi trong bể chứa $3\,200\text{ l}$ nước thì mực nước cao $0{,}8\text{ m}$.}
\loigiai{
\begin{itemchoice}
    \itemch Đúng. Diện tích đáy: $2{,}5 \times 1{,}6 = 4\text{ m}^2$.
    \itemch Đúng. Dung tích tối đa: $4 \times 1{,}2 = 4{,}8\text{ m}^3 = 4\,800\text{ l}$.
    \itemch Sai. Chiều cao mực nước: $3\,200\text{ l} = 3{,}2\text{ m}^3$; $3{,}2 : 4 = 0{,}8\text{ m} \neq 0{,}75\text{ m}$.
    \itemch Đúng. Mực nước cao $0{,}8\text{ m}$.
\end{itemchoice}
}
\end{ex}

\textbf{Phần 3. Trả lời ngắn}
\setcounterpageref{ex}{0}

\begin{ex} %[5H12V71]
Một hình tam giác có độ dài đáy là $35\text{ cm}$, chiều cao tương ứng bằng $\dfrac{4}{5}$ độ dài đáy. Tính diện tích của hình tam giác đó theo xăng-ti-mét vuông.
\shortans{490}
\loigiai{
Chiều cao của tam giác là:
$$35 \times \dfrac{4}{5} = 28\text{ (cm)}$$
Diện tích tam giác là:
$$\dfrac{35 \times 28}{2} = 490\text{ (cm}^2\text{)}$$
\textbf{Đáp số:} $490$
}
\end{ex}

\begin{ex} %[5H12V71]
Một tấm bìa hình tròn có diện tích là $314\text{ cm}^2$. Tính đường kính của tấm bìa hình tròn đó theo xăng-ti-mét (lấy $\pi = 3{,}14$).
\shortans{20}
\loigiai{
Tích của bán kính với bán kính là:
$$314 : 3{,}14 = 100\text{ (cm}^2\text{)}$$
Vì $10 \times 10 = 100$ nên bán kính hình tròn là $10\text{ cm}$.\\
Đường kính của tấm bìa hình tròn là:
$$10 \times 2 = 20\text{ (cm)}$$
\textbf{Đáp số:} $20$
}
\end{ex}

\begin{ex} %[5H12V71]
Một mảnh đất hình thang có đáy lớn $36\text{ m}$, đáy bé $24\text{ m}$ và chiều cao $15\text{ m}$. Người ta dành $20\%$ diện tích mảnh đất để trồng hoa. Tính diện tích đất trồng hoa theo mét vuông.
\shortans{90}
\loigiai{
Diện tích mảnh đất hình thang là:
$$\dfrac{(36 + 24) \times 15}{2} = 450\text{ (m}^2\text{)}$$
Diện tích đất trồng hoa là:
$$450 \times 20 : 100 = 90\text{ (m}^2\text{)}$$
\textbf{Đáp số:} $90$
}
\end{ex}

\begin{ex} %[5H12V71]
Người ta làm một cái thùng không nắp dạng hình lập phương bằng tôn có cạnh dài $8\text{ dm}$. Tính diện tích tôn cần dùng để làm chiếc thùng đó theo đề-xi-mét vuông (không kể mép hàn).
\shortans{320}
\loigiai{
Thùng không nắp gồm $5$ mặt hình vuông.\\
Diện tích tôn dùng làm thùng là:
$$8 \times 8 \times 5 = 320\text{ (dm}^2\text{)}$$
\textbf{Đáp số:} $320$
}
\end{ex}

\begin{ex} %[5H12V71]
Một căn phòng dạng hình hộp chữ nhật có thể tích là $120\text{ m}^3$, chiều dài $6\text{ m}$ và chiều cao $4\text{ m}$. Tính chiều rộng của căn phòng đó theo mét.
\shortans{5}
\loigiai{
Diện tích đáy căn phòng là:
$$120 : 4 = 30\text{ (m}^2\text{)}$$
Chiều rộng của căn phòng là:
$$30 : 6 = 5\text{ (m)}$$
\textbf{Đáp số:} $5$
}
\end{ex}

\begin{ex} %[5H12V71]
Một bể nước hình hộp chữ nhật có chiều dài $18\text{ dm}$, chiều rộng $15\text{ dm}$, chiều cao $12\text{ dm}$. Hiện tại bể đang chứa $1\,620\text{ l}$ nước. Cần đổ thêm bao nhiêu lít nước nữa để đầy bể?
\shortans{1620}
\loigiai{
Thể tích của bể nước là:
$$18 \times 15 \times 12 = 3\,240\text{ (dm}^3\text{)} = 3\,240\text{ (l)}$$
Số lít nước cần đổ thêm để đầy bể là:
$$3\,240 - 1\,620 = 1\,620\text{ (l)}$$
\textbf{Đáp số:} $1620$
}
\end{ex}
